% onepager.tex -- one-page visual for the LinkedIn announcement of the APSEC 2026
% acceptance. 4:5 portrait (20 x 25 cm), the ratio LinkedIn shows without cropping.
% Build: make apsec-onepager  (PDF + PNG at 150 dpi, 1181 x 1476 px).
% Everything here is public: the paper's own abstract, its Figure 1, and the repository.
\documentclass[11pt]{article}

\usepackage[paperwidth=20cm,paperheight=25cm,margin=1.3cm,top=0cm,bottom=1.1cm]{geometry}
\usepackage[T1]{fontenc}
\usepackage[utf8]{inputenc}
\usepackage{helvet}
\renewcommand{\familydefault}{\sfdefault}
\usepackage{graphicx}
\graphicspath{{../figures/}}
\usepackage{xcolor}
\usepackage{microtype}
\usepackage[hidelinks]{hyperref}

\definecolor{brandgreen}{HTML}{00915A}
\definecolor{inkblue}{HTML}{1F3A5F}
\definecolor{softgrey}{HTML}{F2F4F5}

\setlength{\parindent}{0pt}
\pagestyle{empty}

% A flag item: green square, then the text.
\newcommand{\point}[2]{%
  \begin{minipage}[t]{0.028\textwidth}\raisebox{0.15ex}{\textcolor{brandgreen}{\rule{1.6mm}{1.6mm}}}\end{minipage}%
  \begin{minipage}[t]{0.96\textwidth}\textbf{#1} #2\end{minipage}\par\smallskip}

\begin{document}

% --- Venue band ------------------------------------------------------------
\noindent\colorbox{brandgreen}{%
  \begin{minipage}{\dimexpr\textwidth-2\fboxsep\relax}
    \vspace{3mm}
    \textcolor{white}{\large\textbf{Accepted at APSEC 2026}}\hfill
    \textcolor{white}{\normalsize Software Engineering in Practice track}\\[1mm]
    \textcolor{white}{\footnotesize 33rd Asia-Pacific Software Engineering Conference \quad\textbullet\quad Bali, 7--10 December 2026 \quad\textbullet\quad 21 of 57 accepted}
    \vspace{3mm}
  \end{minipage}}

\vspace{5mm}

% --- Title and authors -----------------------------------------------------
{\color{inkblue}\LARGE\textbf{VERA: A Responsible AI Dashboard for Assessing}\\[2mm]
\textbf{LLM Deployments in Light of the EU AI Act}\par}

\vspace{3mm}
{\color{black!65}\footnotesize
Robin Quériaux\textsuperscript{1} \enspace Solveig Bachellery\textsuperscript{1} \enspace Léa Déléris\textsuperscript{1} \enspace Saurabh Rajput\textsuperscript{2} \enspace Mariam Barry\textsuperscript{1} \enspace Tushar Sharma\textsuperscript{2} \enspace Jérôme Lebecq\textsuperscript{1}\\[0.5mm]
\textsuperscript{1}BNP Paribas, France \qquad \textsuperscript{2}Dalhousie University, Canada\par}

\vspace{4mm}

% --- The one-line claim ----------------------------------------------------
{\color{inkblue}\large\textbf{%
The people who approve an AI deployment are not the people who run the benchmarks.\\
VERA gives them evidence they can read, question and file.}\par}

\vspace{4mm}

% --- Figure ----------------------------------------------------------------
\noindent\includegraphics[width=\textwidth]{fig1.png}

\vspace{1mm}
{\color{black!60}\scriptsize One reviewer, one run: pick a connected model, choose the requirements,
run public benchmarks, and read one score per requirement of the EU AI Act --- no command line, no code.\par}

\vspace{4mm}

% --- What it changes --------------------------------------------------------
{\color{inkblue}\large\textbf{What it changes in a review}\par}
\vspace{3mm}
\point{Every score opens.}{A reviewer expands a requirement down to the benchmarks behind it and to the model's own answers, instead of trusting a number on a slide.}
\point{Every artifact names its run.}{Each evaluation emits its own model card, datasheet and export, so a filed document cannot drift from the evaluation it describes.}
\point{Nothing leaves the bank.}{The stack is self-hosted, and the energy each evaluation consumes is measured rather than declared.}

\vspace{3mm}

% --- Abstract ---------------------------------------------------------------
\noindent\colorbox{softgrey}{%
  \begin{minipage}{\dimexpr\textwidth-2\fboxsep\relax}
    \vspace{3mm}
    {\color{inkblue}\textbf{Abstract}}\\[2mm]
    {\scriptsize Institutions that deploy large language models in regulated settings often need
    technical evidence of responsible AI practices before approving a deployment. COMPL-AI provides a
    technical interpretation of the EU AI Act that organises this evidence under six principle groups
    and eighteen technical requirements, twelve of which current benchmarks can measure. The benchmark
    suites that produce this evidence are typically operated from a command line, while the staff who
    approve deployments are non-engineers. We present VERA, an open-source self-hosted dashboard
    deployed to facilitate responsible AI processes within a large European bank. A reviewer selects a
    model and runs public benchmarks against it. VERA scores the twelve requirements that current
    benchmarks can measure and instruments the energy consumed by each evaluation, adding run-level
    evidence for environmental impact. We run four served open-weight models through the requirement
    pipeline, and the dashboard displays a score for each measurable requirement. A within-subject
    study then compares how participants read these benchmark results: from the raw output files in
    one condition, and through the VERA dashboard in the other. The dashboard and the study material
    are released publicly to promote replicability and reproducibility in different organisational
    contexts.\par}
    \vspace{3mm}
  \end{minipage}}

\vfill

% --- Footer -----------------------------------------------------------------
{\color{black!55}\footnotesize
\textbf{Open source:} \href{https://github.com/Rqbln/VERA}{github.com/Rqbln/VERA}
\hfill Authors' accepted version. \copyright~2026 IEEE.\par}

\end{document}
